% =============================================================================
% Lecture 4b: Data Engineering, Part 2: Delivery and Drift
% Introduction to Machine Learning Systems (Volume I)
% =============================================================================
% A released dataset must still reach the accelerator as fast as it computes
% and stay close to the data the world now sends. Running examples: ResNet-50
% on one A100 (delivery); the keyword spotter from Part 1 (drift).
%
% Build: make -C slides build-04b_data_engineering   (figures: images/make_figs.py)
% Plan:  lecture_map.yml and storyboard.md in this folder.
% =============================================================================

\documentclass[aspectratio=169, 12pt]{beamer}
\usepackage{../../assets/beamerthememlsys}
\mlsyssetup{
  volume={Volume I},
  chapter={Chapter 4},
  chaptertitle={Data Engineering, Part 2},
  layer=data,
  logo={../../assets/img/logo-mlsysbook.png},
  instlogo={../../assets/img/logo-harvard.png},
}
\usepackage[T1]{fontenc}
\usepackage[scaled=0.9]{helvet}
\usepackage{courier}
\renewcommand{\familydefault}{\sfdefault}
\usepackage{booktabs,amsmath}
\usetikzlibrary{arrows.meta,positioning}
\graphicspath{{images/}}
\setcounter{mlsystotalsections}{5}

\title{Data Engineering, Part 2: Delivery and Drift}
\author{Vijay Janapa Reddi}
\institute{Harvard University}
\date{}

% A large answer choice for predict and concept-question slides
\newcommand{\choice}[2]{\textbf{\textcolor{mlsysaccent}{#1}}\enspace #2}
% A highlighted (revealed) answer choice
\newcommand{\pick}[2]{\colorbox{mlsysaccent!15}{\choice{#1}{#2}}}
% One stage box in the delivery pipeline
\tikzset{stage/.style={draw=#1, fill=#1!12, line width=1.2pt, rounded corners=3pt,
  minimum width=3.3cm, minimum height=1.05cm, font=\small\bfseries, text=darktext, align=center},
  badge/.style={font=\footnotesize\bfseries, text=#1, align=center}}
% The two silent failures of the hook, drawn once and reused in the synthesis
\newcommand{\hookstall}[1][1]{%
  \begin{minipage}[t]{0.46\textwidth}\centering
    {\footnotesize\bfseries\textcolor{midgray}{TRAINING, NINE VISION AND AUDIO MODELS}}\\[0.25cm]
    {\fontsize{40}{44}\selectfont\bfseries\textcolor{mlsysaccent}{up to 5$\times$}}\\[0.2cm]
    \ifnum#1=1 {\small shorter training from a better data loader:\\same model, same accelerators}\fi
  \end{minipage}}
\newcommand{\hookdrift}[1][1]{%
  \begin{minipage}[t]{0.46\textwidth}\centering
    {\footnotesize\bfseries\textcolor{midgray}{DEPLOYED KEYWORD SPOTTER (ILLUSTRATIVE)}}\\[0.3cm]
    \begin{tikzpicture}[font=\small]
      \draw[line width=1.5pt, vlightgray] (0,0) -- (4.6,0);
      \fill[GreenLine] (0.5,0) circle (0.12); \fill[RedLine] (4.1,0) circle (0.12);
      \node[above=0.12cm, font=\large\bfseries, text=GreenLine] at (0.5,0) {98\%};
      \node[above=0.12cm, font=\large\bfseries, text=RedLine] at (4.1,0) {4 points lower};
      \node[below=0.15cm, align=center] at (0.5,0) {launch};
      \node[below=0.15cm, align=center] at (4.1,0) {six months later};
    \end{tikzpicture}%
    \ifnum#1=1 \\[0.15cm]{\small no one touched the model}\fi
  \end{minipage}}

\begin{document}

% =============================================================================
% TITLE
% =============================================================================
\mlsysnote{[0.5 min] Read the thesis aloud once. Say what today is: Part 1 built a dataset we can trust and released it. Today asks two questions of that release: can it reach the accelerator as fast as the accelerator computes, and does it still look like the data the world sends? ResNet-50 on one A100 carries delivery; the keyword spotter from Part 1 carries drift. The notes total about 71 minutes, which leaves a few minutes for questions. FLEX: [CORE].}
\mlsystitle{Delivery and Drift}{A released dataset must still reach the accelerator as fast as it computes and stay close to the data the world now sends; neither failure raises an error.}{cover_data_engineering_blueprint.png}

% =============================================================================
% HOOK, RECALL, GOALS
% =============================================================================
\begin{frame}{Two systems failed and neither raised an error}
  \mlsysactivity{Discuss}{2 min}
  \centering
  \vspace{0.2cm}
  \hookstall\hfill\uncover<2->{\hookdrift}\\[0.7cm]
  \uncover<3->{\normalsize What number would have exposed each one, and when?}
  \mlsyscredit{Left: Mohan et al., ``Analyzing and Mitigating Data Stalls in DNN Training,'' VLDB 2021. Right: Data Engineering, \S4.12 (illustrative scenario).}
  \note{[3.5 min] NARRATE: Step 1: a measured result. Mohan and colleagues profiled nine widely used vision and audio models and found training time often dominated by data stalls, time the accelerator spends waiting for data to be fetched and preprocessed. A loader that coordinated caching and preprocessing cut training time by up to five times on one server, with the same model and the same accelerators. Nothing crashed; the accelerators were idle much of the time. Step 2: the chapter's illustrative keyword spotter. It launched at 98 percent detection; six months later it measured four points below target, and no one had changed the model. Step 3: the question. ENGAGE: two minutes with a neighbor. Expected for the left: the share of each training step spent waiting on input, set against the compute time. Expected for the right: input statistics compared with the training data, and a measured detection rate on fresh labeled audio. WARN: students often answer ``the loss curve'' or ``accuracy on the test set''; both looked fine in these stories, because the test set is the old data and a stalled run still converges, only slowly. Say that the synthesis at the end returns to both panels with the answer. FLEX: [CORE].}
\end{frame}

\begin{frame}[c]{Recall: what did the dataset release record?}
  \mlsysactivity{Retrieval}{1.5 min}
  \centering
  {\Large\bfseries Close your notes.}\\[0.35cm]
  {\normalsize Name two fields the KWS release records for every clip,\\and the failure each one prevents.}\\[0.4cm]
  \uncover<2->{%
  \footnotesize\color{midgray}
  \begin{tabular}{@{}ll@{}}
    \toprule
    \textbf{Field} & \textbf{Failure it prevents} \\
    \midrule
    Content hash of the audio & duplicates across splits \\
    Split from a hash of the speaker key & one voice in training and test \\
    Label with its rubric version & a rule change mistaken for a data change \\
    Feature configuration and code commit & training-serving skew \\
    \bottomrule
  \end{tabular}}
  \mlsyscredit{Data Engineering, \S4.8 (table 5).}
  \note{[1.5 min] NARRATE: one minute of writing, then take two answers before step 2. Step 2 shows Part 1's release table as the answer key. Expected answers: the split assigned from a hash of the speaker key, so a voice never appears in both training and test; the label with its rubric version, so a change in the labeling rules can be told apart from a change in the data; the feature configuration with its frozen statistics and code commit, so serving computes the same features as training; the content hash, which finds duplicates. LINK: the last segment today uses exactly these fields to trace a regression. FLEX: [CORE]; if short on time, ask it aloud without the writing.}
\end{frame}

\begin{frame}{Today: keep pace with the accelerator and with the world}
  \begin{columns}[c]
    \begin{column}{0.56\textwidth}
      \normalsize
      \begin{itemize}\setlength{\itemsep}{0.3cm}
        \item Price moving a dataset, and pick a storage tier for the access pattern.
        \item Find the stage that binds an input pipeline.
        \item Cut the bytes each sample costs: layout, codec, shards.
        \item Detect drift and trace a regression to its cause.
      \end{itemize}
    \end{column}
    \begin{column}{0.40\textwidth}
      \begin{mlsyscard}{computestroke}
        \textbf{Delivery: ResNet-50, one A100}\\[0.05cm]
        \small 3,200 images/s sustained\\150 KB per JPEG
      \end{mlsyscard}
      \vspace{0.2cm}
      \begin{mlsyscard}{datastroke}
        \textbf{Drift: the keyword spotter}\\[0.05cm]
        \small 98\% detection at launch\\the release from Part 1
      \end{mlsyscard}
    \end{column}
  \end{columns}
  \note{[1.5 min] NARRATE: read the driving question aloud: how do we know the data pipeline is keeping pace with the accelerator and with the world? Four outcomes, one per pair of segments. Two running examples. ResNet-50 image classification on one A100, NVIDIA's datacenter GPU, carries delivery, because images are large enough to stress storage, decoding, and transfer in ways that one-second audio clips do not. The keyword spotter (KWS) from Part 1 carries drift, because its release is what makes a later failure traceable. In D·A·M terms, today is the Data--Machine overlap, and then life after launch. FLEX: [CORE].}
\end{frame}

% =============================================================================
% SECTION: MOVING DATA
% =============================================================================
\mlsysnote{[0.5 min] Data has mass: moving it costs time, money, and energy. Two questions for this segment: what does it cost to move a dataset to the computation, and which storage tier fits which access pattern? FLEX: [CORE].}
\mlsysdivider{What does it cost to move a dataset?}{Data Engineering, \S4.9.1, \S4.10.1, and \S4.10.2}
\section{Moving Data}

\begin{frame}{Moving a petabyte takes about a day and \$20,000}
  \mlsysactivity{Estimate}{3 min}
  \begin{columns}[c]
    \begin{column}{0.30\textwidth}
      \begin{mlsyscard}{mlsysaccent}
        \footnotesize\textbf{Given}\\
        1 PB in US East\\
        Accelerators in US West\\
        Link: 100 Gb/s\\
        Egress: \$0.02 per GB
      \end{mlsyscard}
    \end{column}
    \begin{column}{0.66\textwidth}
      \small
      \uncover<2->{Link: 100 Gb/s $\div$ 8 $= 12.5$ GB/s\\[0.1cm]
      Time: $10^{6}$ GB $/$ 12.5 GB/s $= 80{,}000$ s $\approx$ \textbf{22 h}\\[0.1cm]}
      \uncover<3->{Egress: $10^{6}$ GB $\times$ \$0.02 $=$ \textbf{\$20,000}\\[0.15cm]}
      \uncover<4->{\safeimg[width=\textwidth,height=3.1cm,keepaspectratio]{transfer_time.pdf}\\[0.05cm]
      \textcolor{mlsysaccent}{Large data pulls the computation to it: move the code.}}
    \end{column}
  \end{columns}
  \mlsyscredit{Data Engineering, \S4.9.1 and Napkin Math 4.7 (AWS transfer pricing, 2024).}
  \note{[3.5 min] ENGAGE: give one minute to estimate time and cost before step 2. NARRATE: Step 2: a dedicated 100 gigabit-per-second link moves 12.5 gigabytes per second, so a petabyte takes 80,000 seconds, about 22 hours, at full line rate. Step 3: egress, the fee a cloud provider charges for data leaving its network, at two cents a gigabyte adds \$20,000, about 5,000 hours on Google's TPU v4 accelerator. Step 4: on a 10 Gb/s link the same move takes about 9.3 days; the dashed line is a 20-hour training run. If the run is shorter than its own transfer, or cheaper than \$20,000, move the code to the data. This is data gravity. WARN: the wire time is the optimistic part. The chapter's pitfall is planning a migration as bandwidth alone, when pipelines, validation checks, access controls, and caches must move too. NARRATE: lakehouse engines, which query data where it already sits in object storage, are a direct response to data gravity. LINK: movement also costs energy: the chapter's table 6 puts one 32-bit DRAM access at about 173 times, and one network transfer at about 10,811 times, the energy of an FP32 multiply; Lecture 11 develops this. Removing redundant samples before they move is the other answer, and Lecture 9 develops data selection. FLEX: [CORE].}
\end{frame}

\begin{frame}{Storage got 10 orders of magnitude cheaper; the tier gaps stayed}
  \centering
  \mlsyssteps[width=0.92\textwidth,height=5.3cm,keepaspectratio]{price_trend_step}{2}
  \mlsyscredit{Data Engineering, \S4.10.2 (figure 16, redrawn); data from Our World in Data, after McCallum.}
  \note{[2.5 min] NARRATE: DRAM is the server's main memory; solid-state drives and hard disks hold data that outlives a power cycle. The plot shows the lowest recorded price per terabyte of each medium, in constant 2020 dollars, on a log axis; every gridline is a factor of 100,000. Step 1: the curves. Every medium fell. Step 2: hard disk went from about \$88 billion per terabyte in 1956 to \$10.6 in 2023, roughly ten orders of magnitude. Yet in 2023 DRAM still cost 42 times as much as solid-state storage per terabyte, and solid-state storage 2.4 times as much as hard disk. Since the 1980s the gaps between the curves have not closed. ENGAGE: if capacity is nearly free, why does a team still choose tiers? Expected: because no single tier is both cheap and fast; the choice of tier, not the capacity, sets the cost of keeping a training job fed. WARN: students read ``storage is cheap'' as ``storage does not matter''; cheap capacity and fast delivery are different purchases. FLEX: [OPTIONAL]; say the last line aloud if cut.}
\end{frame}

\begin{frame}{Serving needs low latency; training needs high bandwidth}
  \centering
  \begin{columns}[T]
    \begin{column}{0.49\textwidth}
      \centering\safeimg[width=\textwidth,height=4.4cm,keepaspectratio]{latency_ladder.pdf}\\[0.05cm]
      \uncover<3->{\small\textbf{\textcolor{BlueLine}{Serving:}} 10,000 requests/s $\times$ 10 reads\\$=$ 100,000 random reads per second}
    \end{column}
    \begin{column}{0.49\textwidth}
      \centering\uncover<2->{\safeimg[width=\textwidth,height=4.4cm,keepaspectratio]{tier_bandwidth.pdf}\\[0.05cm]}
      \uncover<3->{\small\textbf{\textcolor{OrangeLine}{Training:}} one long sequential scan\\of every image, every epoch}
    \end{column}
  \end{columns}
  \mlsyscredit{Data Engineering, \S4.10 (margin figure), \S4.10.1, and \S4.10.2 (table 8).}
  \note{[4 min] Define the tiers: DRAM is main memory; local NVMe is flash inside the server; a remote SSD is flash over the network; an object store such as S3 serves whole files over the web. NARRATE: Step 1: latency, log scale; the dashed line is 1 ms, a typical budget for one feature lookup. Only DRAM (100 ns) and local NVMe (100 µs) fit. Step 2: sequential bandwidth from table 8; note the object store, 100 to 500 MB/s per connection. Step 3: a recommendation service at 10,000 requests a second, ten feature reads each, needs 100,000 random reads a second (IOPS): a latency problem. A feature store splits along this line: an online key-value store serves those reads, and an offline store holds the same feature definitions in columnar files for training. A training job scans every image in order each epoch: a bandwidth problem. Say aloud: the KWS corpus, 748.8 GB, costs \$17.2 a month in object storage and loads in 7,488 s; on NVMe, \$74.9 to \$224.6 a month and 149.8 s. ENGAGE: where do the raw audio and this week's working set live? Expected: object store; NVMe. Landing line: each tier fits an access pattern, latency to serve and bandwidth to train. FLEX: [CORE].}
\end{frame}

% =============================================================================
% SECTION: FEEDING
% =============================================================================
\mlsysnote{[0.5 min] From moving a dataset once to delivering it every epoch. ResNet-50 on one A100 is the running example. FLEX: [CORE].}
\mlsysdivider{Can storage keep the accelerator fed?}{Data Engineering, \S4.9.3, \S4.10.3, and \S4.10.4}
\section{Feeding}

\begin{frame}{The A100 can consume 12,682 images a second; it sustains 3,200}
  \centering
  \begin{tikzpicture}[>={Stealth[length=3mm]}, line width=1.2pt]
    \node[stage=midgray] (s) at (0,0) {Storage\\[-1pt]\footnotesize\mdseries JPEG files};
    \node[stage=BlueLine] (h) at (4.8,0) {Host CPU\\[-1pt]\footnotesize\mdseries decode, augment};
    \node[stage=OrangeLine] (a) at (9.6,0) {Accelerator\\[-1pt]\footnotesize\mdseries train};
    \draw[->, midgray] (s) -- (h); \draw[->, midgray] (h) -- (a);
    \node[badge=OrangeLine, below=0.2cm of a] {3,200 images/s};
    \uncover<2->{\node[badge=BlueLine, below=0.2cm of h] {decode is CPU work};}
    \uncover<3->{\node[badge=RedLine, below=0.2cm of s] {must supply ? MB/s};}
  \end{tikzpicture}\\[0.35cm]
  \safeimg[width=0.80\textwidth,height=2.4cm,keepaspectratio]{appetite.pdf}
  \mlsyscredit{Data Engineering, \S4.9.3: 312 TFLOP/s $\div$ 24.6 GFLOP per training image $=$ 12,682 images/s.}
  \note{[2.5 min] NARRATE: one diagram carries this segment and the next: storage, the host CPU, and the accelerator, each a stage the data passes through every epoch. Step 1: the accelerator. Its peak arithmetic rate is 312 TFLOP/s, trillion floating-point operations a second; one training step on one image costs about 24.6 GFLOP, roughly three times a prediction because training also computes how to adjust every parameter. The arithmetic ceiling is therefore 12,682 images a second. A well-tuned pipeline sustains about 3,200, roughly 25 percent of the ceiling, once layers that wait on memory and synchronization take their share. Step 2: the host CPU reads each JPEG, decodes it into pixels, and applies random crops and flips (augmentation) before handing a batch to the accelerator. Step 3: storage must deliver the bytes. How many megabytes a second? That is the next slide. WARN: students assume the accelerator is always the bottleneck because it is the expensive part. The overlapped form of the iron law from Lecture 1 says a pipeline runs at the rate of its slowest stage, whichever that is. FLEX: [CORE].}
\end{frame}

\begin{frame}{Your turn: can 250 MB/s keep ResNet-50 fed?}
  \mlsysactivity{Estimate}{3 min}
  \begin{columns}[T]
    \begin{column}{0.40\textwidth}
      \begin{mlsyscard}{mlsysaccent}
        \textbf{Given}\\[0.1cm]
        One A100 trains ResNet-50 at 3,200 images/s\\
        Each image: 150 KB JPEG\\
        Storage supplies 250 MB/s
      \end{mlsyscard}
    \end{column}
    \begin{column}{0.56\textwidth}
      \normalsize
      \uncover<2->{Needed: $3{,}200 \times 150$ KB $=$ \textbf{480 MB/s}\\[0.3cm]}
      \uncover<3->{Supplied: $250 / 480 = 0.52$\\[0.3cm]}
      \uncover<4->{{\Large\textbf{\textcolor{mlsysaccent}{48\% of the training rate lost}}}}
    \end{column}
  \end{columns}
  \mlsyscredit{Data Engineering, \S4.9.3.}
  \note{[3 min] ENGAGE: one minute with a neighbor before step 2. Step 2: the accelerator consumes 3,200 images a second; at 150 KB each, storage must deliver 480 MB/s. Step 3: an attached cloud block volume delivers 250 MB/s, so it can sustain 52 percent of the training rate. Step 4: the modeled feeding loss is 48 percent: the accelerator waits on input for about half of every step. If many answer 52 percent lost, they confused the surviving fraction with the loss. WARN: this is a loss relative to the sustained 3,200 images a second, not a measurement of all idle time against the advertised peak; the chapter is careful about that distinction, and the slide keeps that distinction. NARRATE: this is the stall that the hook's first panel measured across nine real models. FLEX: [CORE].}
\end{frame}

\mlsysnote{[1 min] The segment's law, the chapter's equations 4 and 5. Training throughput is bounded by the smaller of two rates: what the accelerator can compute, and what the input path can supply. The supply rate is storage bandwidth, less overhead, divided by the bytes each sample costs. Once supply is the smaller rate, a faster accelerator buys nothing; the two levers are more bandwidth or fewer bytes per sample. The next segment is about the second lever. LINK: this is the data term of the iron law from Lecture 1, in its overlapped form: a pipeline runs at the rate of its slowest stage. FLEX: [CORE].}
\mlsysfocus{Training runs at the slower of compute and data supply}{%
  $\displaystyle \text{Throughput} \;\le\; \min\Big(\underbrace{R_{\text{compute}}}_{\text{accelerator}},\;
   \underbrace{\frac{BW_{\text{storage}}\,(1-\text{overhead})}{\text{bytes per sample}}}_{\text{data supply}}\Big)$\\[0.6cm]
  {\small\textcolor{white!75}{Samples per second on both sides; equations 4 and 5 of the chapter.}}%
}

\begin{frame}{Caching decoded tensors makes the loss worse}
  \mlsysactivity{Predict}{2 min}
  \centering
  {\normalsize To spare the CPU, the team stores decoded FP32 tensors instead of JPEGs.}\\[0.25cm]
  {\large\choice{A}{faster}\qquad\choice{B}{no change}\qquad\alt<2->{\pick{C}{slower}}{\choice{C}{slower}}}\\[0.35cm]
  \uncover<2->{\mlsyssteps[width=0.90\textwidth,height=3.1cm,keepaspectratio]{decoded_bytes_step}{2}\\[0.1cm]}
  \uncover<3->{\small 602 KB per image, 4$\times$ the JPEG. Fixed tensors also freeze the random crops.}
  \mlsyscredit{Data Engineering, \S4.9.3.}
  \note{[3 min] ENGAGE: vote with fingers, then reveal. Step 2: the bars. Storage still supplies 250 MB/s (dashed line). The JPEG path needs 480 MB/s; a decoded image stored as 32-bit floating-point numbers (FP32) is 602 KB, about four times the compressed JPEG, so the same 3,200 images a second now need 1.9 GB/s, and the modeled loss rises from 48 to 87 percent. Step 3: a second cost: augmentation applies a different random crop to each image every epoch; tensors fixed on disk freeze one crop forever. If many chose A, they reasoned from the CPU alone; the bytes per sample went up fourfold, and storage was already the slowest stage. WARN: ``which bytes count'' matters as much as how many. Compression trades storage bandwidth for CPU decode work, and the right side of that trade depends on which stage binds. NARRATE: many loaders ship 8-bit pixels and convert to floating point on the accelerator for this reason. FLEX: [CORE].}
\end{frame}

\begin{frame}{A SATA SSD feeds today's rate, not the ceiling or eight A100s}
  \begin{columns}[c]
    \begin{column}{0.76\textwidth}
      \centering\mlsyssteps[width=\textwidth,height=5.0cm,keepaspectratio]{storage_supply_step}{3}
    \end{column}
    \begin{column}{0.20\textwidth}
      \centering\safeimg[width=\textwidth,height=3.6cm,keepaspectratio]{nvme-ssd-1tb.jpg}\\[0.1cm]
      {\footnotesize\textcolor{midgray}{An NVMe SSD}}
    \end{column}
  \end{columns}
  \mlsyscredit{Data Engineering, Napkin Math 4.8 and \S4.10.1. Photo: User5515, CC0, Wikimedia Commons.}
  \note{[2.5 min] NARRATE: Step 1: the orange bar is what ResNet-50 needs today, 480 MB/s. One object-store connection streams about 100 MB/s, so even today's rate needs parallel streams of large shards. The attached block volume, 250 MB/s, is the 48 percent loss we just computed. A SATA SSD at about 500 MB/s supplies 3,333 images a second, just enough. Step 2: the compute ceiling, 12,682 images a second, needs 1.9 GB/s of compressed JPEGs. The object store would need 19 parallel streams; the SATA drive holds the accelerator to 26 percent of the ceiling; only local NVMe at 5 to 7 GB/s covers it. Step 3: eight A100s at the sustained rate need 25,600 images a second, 3.8 GB/s compressed, and 15.4 GB/s if the team stages decoded tensors. ENGAGE: why buy the NVMe if SATA suffices today? Expected: each fix to the slowest stage exposes the next; size storage for the rate the pipeline is heading toward. WARN: a dataset that fits in host memory, as ImageNet does on many eight-accelerator servers, is read from storage once and then served from the page cache, the copy of recently read files the operating system keeps in memory, so this budget binds on datasets larger than memory, shared nodes, and remote streaming. FLEX: [CORE].}
\end{frame}

\begin{frame}{With storage fixed, host decode binds next}
  \centering
  \safeimg[width=0.86\textwidth,height=4.6cm,keepaspectratio]{dataloader.pdf}\\[0.15cm]
  \uncover<2->{\small Nine vision and audio models: data stalls often dominate training time (Mohan et al.\ 2021).}
  \mlsyscredit{Data Engineering, \S4.10.4 (figure 17, redrawn; the dashed loader capacity stops at the figure's illustrative 3,500 images/s cap).}
  \note{[2 min] NARRATE: assume storage now supplies 480 MB/s. The host CPU must still decode and augment 3,200 images a second. Each loader worker process handles about 375 images a second, so throughput rises with workers until about nine workers match the accelerator. Left of the crossing, the host binds and adding accelerators buys nothing; right of it, the accelerator binds and more workers buy nothing. Above the crossing the loader could supply more, the thin dashed line, but the accelerator takes only 3,200. For keyword spotting, MFCC extraction and audio augmentation load the same host threads. Step 2: the measured version. Mohan and colleagues found that once storage is provisioned, CPU-side preparation is the common stall. ENGAGE: eight A100s on one host need how many decode workers? Expected: about 72, more cores than many hosts can spare, which is why hardware JPEG decoders and prefetching exist. Landing line for the segment: training runs at the slower of compute and data supply, and each fix to the slowest stage exposes the next. LINK: Lecture 8 counts the workers for an H100 and overlaps loading with compute; here we only see where the bottleneck moves. FLEX: [OPTIONAL]; say the landing line if cut.}
\end{frame}

% =============================================================================
% SECTION: FEWER BYTES
% =============================================================================
\mlsysnote{[0.5 min] The law has two levers. We have tried more bandwidth; now the other one, fewer bytes per sample. FLEX: [CORE].}
\mlsysdivider{How do we cut the bytes each sample costs?}{Data Engineering, \S4.10.5 and \S4.11}
\section{Fewer Bytes}

\begin{frame}{Columnar files read one fifth of the bytes; dictionaries cut more}
  \centering
  \mlsyssteps[width=0.84\textwidth,height=2.6cm,keepaspectratio]{format_bytes_step}{2}\\[0.3cm]
  \uncover<3->{%
  \begin{tikzpicture}[>={Stealth[length=2.5mm]}, font=\small, every node/.style={align=center}]
    \node[draw=midgray, fill=midgray!10, rounded corners=3pt, minimum width=4.2cm, minimum height=1.1cm] (a) at (0,0) {country\_code, 100M rows\\200 distinct values};
    \node[draw=GreenLine, fill=GreenL, rounded corners=3pt, minimum width=5.2cm, minimum height=1.1cm] (b) at (7.0,0) {dictionary of 200 strings\\+ one small index per row};
    \draw[->, line width=1.2pt, GreenLine] (a) -- node[above, font=\footnotesize\bfseries] {20 to 50$\times$} (b);
  \end{tikzpicture}\\[0.2cm]
  {\small Projection and compression together: 20 to 100$\times$ fewer bytes than uncompressed rows.}}
  \mlsyscredit{Data Engineering, Napkin Math 4.9 (20 of 100 equal-size columns) and \S4.10.5.}
  \note{[2 min] NARRATE: a fraud model uses 20 of a table's 100 columns, all of equal size. Step 1: a row format such as CSV stores each record together, so reading 20 columns means reading all 100. Step 2: a columnar format such as Parquet stores each column together, so the loader reads only those 20, one fifth of the bytes: the same gain as a five-times faster disk (equation 6, effective bandwidth equals physical bandwidth times the useful fraction). Step 3: within a column, values repeat. A country code with 200 distinct values across 100 million rows compresses 20 to 50 times by storing each string once and a small integer per row, dictionary encoding. Projection and compression together cut I/O 20 to 100 times against uncompressed rows. ENGAGE: why does a columnar format not help a model that uses every column? Expected: projection saves nothing; only compression remains. WARN: text formats also cost CPU: parsing CSV or JSON spends hundreds of cycles per sample on string conversion, the chapter's serialization pitfall. FLEX: [OPTIONAL].}
\end{frame}

\begin{frame}{The right codec depends on which stage binds}
  \mlsysactivity{Diagnose}{3 min}
  \begin{columns}[T]
    \begin{column}{0.40\textwidth}
      \begin{mlsyscard}{mlsysaccent}
        \small\textbf{Two loaders, same 100 GB}\\[0.1cm]
        gzip: 6--8$\times$ smaller, 120 MB/s\\
        Snappy: 2--3$\times$, 500 MB/s\\[0.1cm]
        Loader 1: CPU 100\%, disk idle\\
        Loader 2: disk saturated, CPU idle\\[0.1cm]
        Which codec for each?
      \end{mlsyscard}
    \end{column}
    \begin{column}{0.57\textwidth}
      \centering
      \uncover<2->{\safeimg[width=0.9\textwidth,height=2.6cm,keepaspectratio]{codec_minutes.pdf}\\[0.05cm]
      {\small 50 epochs: about 9 h of host CPU difference}\\[0.3cm]}
      \uncover<3->{\small\raggedright
      \textbf{Loader 1} (CPU-bound): Snappy.\\[0.05cm]
      \textbf{Loader 2} (storage-bound): the denser codec;\\fewer \alert{bytes per sample} raise the supply rate.}
    \end{column}
  \end{columns}
  \mlsyscredit{Data Engineering, \S4.10.5.}
  \note{[3.5 min] ENGAGE: pairs decide a codec for each loader before step 2. NARRATE: Step 2: the chapter's numbers. gzip compresses 6 to 8 times but decompresses at about 120 MB/s; Snappy compresses only 2 to 3 times but decompresses at about 500 MB/s. For 100 GB per epoch, gzip costs 13.9 minutes of CPU decompression and Snappy 3.3; over 50 epochs the difference is about 9 hours of host CPU (8.8 hours computed; the chapter rounds to 9). Step 3: the answer depends on which stage binds. Loader 1 is CPU-bound, so CPU minutes are on the critical path: choose Snappy. Loader 2 is storage-bound, so the denser codec wins, because fewer bytes per sample raise the data supply rate in equation 5; the extra CPU work hides behind the disk. WARN: students want one best codec. There is none; there is the codec that relieves the binding stage. If decompression is off the critical path, with spare cores and prefetching, it costs nothing. FLEX: [CORE].}
\end{frame}

\begin{frame}{Shards turn seeks into streams, at the price of shuffling}
  \begin{columns}[c]
    \begin{column}{0.54\textwidth}
      \centering
      \begin{tikzpicture}[font=\footnotesize, every node/.style={align=center}]
        \foreach \i in {0,...,7} \foreach \j in {0,...,3}
          \fill[midgray!60] (\i*0.28,\j*0.28) rectangle ++(0.22,0.22);
        \node[below] at (1.05,-0.05) {millions of 150 KB files:\\one seek each};
        \draw[->, line width=1.2pt, >={Stealth[length=2.5mm]}] (2.5,0.5) -- (3.3,0.5);
        \foreach \j in {0,...,3} \fill[GreenLine] (3.5,\j*0.28) rectangle ++(2.4,0.22);
        \node[below] at (4.7,-0.05) {a few large shards:\\read in sequence};
        \uncover<2->{
          \draw[BlueLine, line width=1.2pt, rounded corners=3pt] (3.4,-2.15) rectangle (6.0,-1.5);
          \node[text=BlueLine, font=\footnotesize\bfseries] at (4.7,-1.825) {shuffle buffer};
          \node[text=darktext, below] at (4.7,-2.15) {mixes records before each batch};}
      \end{tikzpicture}
    \end{column}
    \begin{column}{0.44\textwidth}
      \centering\uncover<3->{\safeimg[width=\textwidth,height=2.6cm,keepaspectratio]{stragglers.pdf}}
    \end{column}
  \end{columns}
  \vspace{0.3cm}
  \uncover<4->{\centering\small FFCV: ResNet-50 to 75\% on ImageNet in about 20 minutes on one machine.\par}
  \mlsyscredit{Data Engineering, \S4.10.5 and \S4.11; Leclerc et al., ``FFCV,'' CVPR 2023.}
  \note{[3.5 min] NARRATE: Step 1: opening millions of 150 KB JPEGs spends more time on seeks than on moving bytes. Shard formats such as WebDataset or TFRecord pack records into a few large files read front to back. Step 2: the price. Records inside a shard arrive in a fixed order, so the loader shuffles the list of shards and mixes records in a buffer before forming each batch; too small a buffer feeds the model correlated batches. Step 3: partitions must be balanced. In an eight-accelerator job, seven workers load their batch in 12 ms; one reads an oversized partition and takes 180 ms, and every accelerator waits for it, because the accelerators exchange gradients together at the end of each step. Step 4: redesigning format and loader together pays: FFCV trains ResNet-50 on ImageNet to 75 percent accuracy in about 20 minutes on one machine. ENGAGE: does the ResNet-50 pipeline now pass its acceptance test on a SATA drive? Expected: yes, today, if images are read as large sequential shards (3,333 against 3,200 images a second); it fails once the rate rises toward the ceiling. Landing line: layout, codec, and shards change bytes per sample, the other lever in the supply rate, and the right choice depends on which stage binds. FLEX: [OPTIONAL]; keep the landing line.}
\end{frame}

% =============================================================================
% SECTION: DRIFT
% =============================================================================
\mlsysnote{[0.5 min] The pipeline is fast enough. Now the second question: does the data the world sends still look like the data we trained on? FLEX: [CORE].}
\mlsysdivider{How do we detect that the world has moved?}{Data Engineering, \S4.5.3, \S4.6.4, and \S4.12.1}
\section{Drift}

\begin{frame}{Re-collecting a test set the same way cost 11 to 14 points}
  \centering
  \safeimg[width=0.90\textwidth,height=3.5cm,keepaspectratio]{recht_ranges.pdf}\\[0.4cm]
  \uncover<2->{\normalsize Same collection procedure, new photos, years later.}
  \mlsyscredit{Recht, Roelofs, Schmidt, and Shankar, ``Do ImageNet Classifiers Generalize to ImageNet?'' ICML 2019 (abstract); cited in Data Engineering, \S4.4.2.}
  \note{[1.5 min] NARRATE: Recht and colleagues rebuilt the CIFAR-10 and ImageNet test sets by following the original collection procedure as closely as they could, with new images. Every model they tested lost accuracy: 3 to 15 points on CIFAR-10 and 11 to 14 points on ImageNet. Step 2: nothing about the task changed; the new photos came from a slightly different moment and source pool. ENGAGE: if a careful re-collection costs this much, what happens to a deployed model whose inputs change every month? Expected: it loses accuracy without any error, unless something measures the shift. WARN: students guess that years of tuning on the old test set caused the drop; the authors found little sign of that, and attribute it to small, natural differences between the old and new images. FLEX: [OPTIONAL]; one sentence aloud if cut.}
\end{frame}

\begin{frame}{Drift is a change in $p(x)$, $p(y)$, or $p(y \mid x)$}
  \centering
  \mlsyssteps[width=0.94\textwidth,height=5.6cm,keepaspectratio]{drift_step}{3}
  \mlsyscredit{Data Engineering, \S4.12.1.}
  \note{[2 min] NARRATE: drift is the world moving away from the training data. Read the notation aloud: $p(x)$ is how often each input occurs, $p(y)$ how often each answer occurs, and $p(y \mid x)$ the chance of each answer given an input, the rule the model learned. Step 1: covariate shift: inputs change, the input-to-label rule does not. The KWS case: a new microphone changes the audio's frequency response; people say the same words. Detect it by comparing input feature distributions. Step 2: label shift: the mix of outcomes changes, such as more positive flu tests in winter, while each class looks the same. Detect it from the model's prediction mix, before labels arrive. Step 3: concept drift: the rule itself changes, as when fraudsters adapt. Only fresh labels reveal it, and labels arrive late. ENGAGE: which is hardest to detect, and why? Expected: concept drift, because the inputs can look unchanged. WARN: drift is not training-serving skew. Skew comes from inside our pipeline (Part 1); drift comes from the world. FLEX: [CORE].}
\end{frame}

\begin{frame}{Vote: which kind of shift is it?}
  \mlsysactivity{Vote}{2.5 min}
  \begin{columns}[c]
    \begin{column}{0.66\textwidth}
      \normalsize
      \begin{tabular}{@{}P{0.70\textwidth}l@{}}
        \toprule
        \textbf{Symptom} & \textbf{Shift} \\
        \midrule
        A new speaker model ships with a different microphone. & \uncover<2->{\textbf{\textcolor{BlueLine}{Covariate}}} \\ \addlinespace[0.25cm]
        In winter, twice as many tests come back positive. & \uncover<3->{\textbf{\textcolor{OrangeLine}{Label}}} \\ \addlinespace[0.25cm]
        Fraudsters learn which purchases the model ignores. & \uncover<4->{\textbf{\textcolor{RedLine}{Concept}}} \\
        \bottomrule
      \end{tabular}
    \end{column}
    \begin{column}{0.30\textwidth}
      \centering\safeimg[width=\textwidth,height=3.8cm,keepaspectratio]{echo-dot-mic-board.jpg}\\[0.1cm]
      {\footnotesize\textcolor{midgray}{Echo Dot microphone board}}
    \end{column}
  \end{columns}
  \mlsyscredit{Photo: Raimond Spekking, CC BY-SA 4.0, Wikimedia Commons.}
  \note{[2.5 min] ENGAGE: students write C, L, or K for each row, compare with a neighbor, then reveal one row per step. Row 1, covariate: the photo shows the ring of microphones on a smart speaker's board; a different microphone changes the audio, not what words mean. Row 2, label: the prevalence of the outcome changed; symptoms still predict disease the same way. Row 3, concept: the same purchase features now mean something different. Expect debate on row 3: some will say the inputs changed too; both can happen at once, and the test is whether the input-to-label mapping moved. WARN: students who answer ``concept'' for row 1 think any change in what the model sees is a change in meaning; ask whether a person listening would hear a different word. FLEX: [CORE].}
\end{frame}

\begin{frame}{The K-S test flags drift when the CDFs separate}
  \begin{columns}[c]
    \begin{column}{0.60\textwidth}
      \centering
      \mlsyssteps[width=\textwidth,height=4.9cm,keepaspectratio]{ks_cdf_step}{3}
    \end{column}
    \begin{column}{0.38\textwidth}
      \normalsize
      \uncover<3->{$D = \max_x \lvert F_{\text{train}}(x) - F_{\text{recent}}(x)\rvert$\\[0.3cm]
        $\displaystyle D_{\text{crit}} \approx 1.36\sqrt{\frac{n_0+n_t}{n_0\, n_t}}$\\[0.2cm]
        $n_0 = n_t = 1{,}000$: $0.061$\\[0.3cm]
        $0.08 > 0.061$: \textbf{\textcolor{mlsysaccent}{flag}}}
    \end{column}
  \end{columns}
  \mlsyscredit{Data Engineering, \S4.5.3 and Napkin Math 4.5. The curves are illustrative; the numbers are the chapter's.}
  \note{[2 min] NARRATE: the scenario from \S4.5.3: a phone app update halves the microphone gain on 15 percent of a batch of crowdsourced clips. Every file has the right format, length, and sampling rate, so the schema check from Part 1 passes. Step 1: the empirical cumulative distribution function (CDF) of training loudness: at each loudness, the share of clips at most that loud. Loudness is in dBFS, decibels relative to the loudest value the recorder can store, so values nearer zero are louder. Step 2: the CDF of 1,000 recent clips, slightly quieter. Step 3: the Kolmogorov-Smirnov statistic $D$ is the largest vertical gap between the curves, here 0.08. For two samples of 1,000, the critical value at the 5 percent significance level is $1.36\sqrt{2000/10^6} = 0.061$; 0.08 exceeds it, so the test flags the feature. ENGAGE: does a flag mean retrain? Expected: no; it starts an investigation and an outcome check. WARN: K-S is for one continuous feature; categorical features need frequency comparisons. FLEX: [CORE].}
\end{frame}

\begin{frame}{At 100,000 clips, harmless shifts flag too}
  \centering
  \safeimg[width=0.80\textwidth,height=4.9cm,keepaspectratio]{ks_vs_n.pdf}\\[0.2cm]
  {\small The critical value falls as $1/\sqrt{n}$: 0.061 at 1,000 clips, 0.006 at 100,000.}
  \mlsyscredit{Data Engineering, Napkin Math 4.5; the value at 100,000 uses the same formula.}
  \note{[1 min] NARRATE: the same formula, plotted against sample size. The critical value shrinks as one over the square root of the number of clips. At 1,000 clips per sample it is 0.061, and the observed 0.08 sits above it. At 100,000 clips per sample the bar falls to about 0.006, so a monitor that sees all traffic flags shifts too small to matter. Large monitors therefore need a threshold on the size of the shift, not only on significance. WARN: students read a significant result as an important one; at this scale significance is nearly free. FLEX: [CORE].}
\end{frame}

\begin{frame}{PSI adds every bin that moved: 0.135 here}
  \centering
  {\small $\text{PSI} = \sum_{\text{bins}} \big(\text{recent} - \text{training}\big)\,\ln\dfrac{\text{recent}}{\text{training}}$}\\[0.1cm]
  \mlsyssteps[width=0.84\textwidth,height=4.7cm,keepaspectratio]{psi_bins_step}{3}
  \mlsyscredit{Data Engineering, \S4.12.1 (formula and bands). Bin shares are illustrative.}
  \note{[3.5 min] NARRATE: the population stability index (PSI) bins a feature, here clip loudness, by the training data's quintiles, so each bin holds 20 percent of training clips (blue). Step 1: the same quieter-microphone shift, binned: recent clips (red) pile into the quiet bins, 30 and 25 percent, and thin out of the loud ones, 15 and 10. These shares are illustrative, chosen for the slide; the formula and bands are the chapter's. Step 2: each bin's term is the difference in share times the log of the ratio. Both factors share a sign, so every bin that moved adds to the index: the quietest bin adds $0.10 \times \ln 1.5 = 0.041$, the loudest $0.10 \times \ln 2 = 0.069$. Step 3: the sum is 0.135. Practitioners treat below 0.1 as stable, 0.1 to 0.25 as moderate, above 0.25 as a prompt to investigate; these are conventions to calibrate against your own baseline, not physical limits. ENGAGE: what would PSI be if the recent shares matched training exactly? Expected: zero. NARRATE: KL divergence is the related measure; it is asymmetric, and it becomes infinite where the recent data fills a bin the baseline left empty, so monitors fix its direction and smooth empty bins. LINK: Lecture 14b puts K-S and PSI on a production dashboard. FLEX: [CORE].}
\end{frame}

\begin{frame}{Input monitors miss three failures}
  \mlsysactivity{Vote}{2 min}
  \centering
  \small
  \renewcommand{\arraystretch}{1.3}
  \begin{tabular}{@{}P{0.27\textwidth}P{0.32\textwidth}P{0.31\textwidth}@{}}
    \toprule
    \textbf{Failure} & \textbf{Why input statistics miss it} & \textbf{What does see it} \\
    \midrule
    Missed wake words & a miss uploads no recording & consented recording studies \\
    \uncover<2->{Label quality drift} & \uncover<2->{$\kappa$ fell 0.85 $\rightarrow$ 0.72; features flat} & \uncover<2->{tracking annotator agreement} \\
    \uncover<3->{Concept drift} & \uncover<3->{inputs can look unchanged} & \uncover<3->{outcome checks on fresh labels} \\
    \bottomrule
  \end{tabular}\\[0.4cm]
  \uncover<4->{\normalsize True or false: \emph{production logs show where the model fails}. \textbf{\textcolor{RedLine}{False.}}}
  \mlsyscredit{Data Engineering, \S4.6.4 (illustrative campaign), \S4.12.1, and \S4.13.}
  \note{[3 min] NARRATE: every monitor so far compares inputs. Three failures slip past them. Row 1: the deployed detector uploads audio, with consent, only when it wakes, so the clips the team sees are correct detections and false wakes. A wake word it misses leaves no recording; the user simply says it again. Measuring misses takes a separate, consented collection. Row 2: from Part 1, new annotators without the rubric training accepted near-homophones; Cohen's kappa, agreement beyond chance, fell from 0.85 to 0.72 in six months while every audio feature stayed flat. Row 3: concept drift needs labels, and labels arrive late. Step 4: ENGAGE: thumbs up or down, then reveal. Two more to put to the room: ``a drift flag means retrain'' (false: it starts an investigation and an outcome check) and ``flat input statistics mean the labels are fine'' (false: row 2). WARN: retraining on the logs alone over-represents the cases the model already handles. Landing line for the segment: input statistics flag covariate shift cheaply, label shift shows in the prediction mix, and concept drift and missed detections need labels and collection designed for them. FLEX: [CORE].}
\end{frame}

% =============================================================================
% SECTION: DEBT
% =============================================================================
\mlsysnote{[0.5 min] Back to the hook's keyword spotter. A monitor flagged it; now trace the cause. FLEX: [CORE].}
\mlsysdivider{Why does a pipeline decay when no one touches it?}{Data Engineering, \S4.12.2 to \S4.12.5}
\section{Debt}

\begin{frame}{Debug a data regression by comparing fixed versions}
  \mlsysactivity{Diagnose}{3 min}
  \centering
  {\small Detection fell four points in six months. A new microphone shipped, the labeling\\rubric was revised, and the metadata schema changed. Which comparison first?}\\[0.15cm]
  \uncover<2->{%
  \small
  \renewcommand{\arraystretch}{1.0}
  \begin{tabular}{@{}P{0.22\textwidth}P{0.36\textwidth}P{0.34\textwidth}@{}}
    \toprule
    \textbf{Old vs.\ new} & \textbf{Hold fixed} & \textbf{A change reveals} \\
    \midrule
    Labeling rubric & predictions on fixed clips & the definition of correct moved \\
    Preprocessing & raw recordings, weights, fitted state & a processing-path change \\
    Microphones & one pipeline, one rubric & a hardware coverage gap \\
    \bottomrule
  \end{tabular}}\\[0.15cm]
  \uncover<3->{%
  \begin{tikzpicture}[font=\footnotesize, >={Stealth[length=2mm]},
      chk/.style={draw=BlueLine, fill=BlueLine!10, rounded corners=3pt, line width=1pt,
        minimum width=2.9cm, minimum height=0.8cm, inner sep=1.5pt, align=center}]
    \node[chk] (c1) at (0,0) {Drift over\\time?};
    \node[chk] (c2) at (3.55,0) {Train vs.\\validation?};
    \node[chk] (c3) at (7.1,0) {Validation vs.\\production?};
    \node[chk] (c4) at (10.65,0) {Subgroup\\gap?};
    \draw[->, line width=1pt, midgray] (c1) -- (c2);
    \draw[->, line width=1pt, midgray] (c2) -- (c3);
    \draw[->, line width=1pt, midgray] (c3) -- (c4);
  \end{tikzpicture}}\\[0.1cm]
  \uncover<4->{\small\textcolor{mlsysaccent}{Each comparison needs the release that recorded what was built.}}
  \mlsyscredit{Data Engineering, \S4.12 (illustrative scenario), \S4.12.5, and figure 18 (order of checks, redrawn).}
  \note{[4 min] ENGAGE: pairs pick a first comparison and say what they would hold fixed, one minute. Step 2: the three controlled comparisons. Change one thing, hold the rest fixed. Row 1 is usually first because it needs no new audio: rescore the same fixed predictions under the old and revised rubrics; a changed score means the definition of correct moved, not the detector. Row 2: replay the same raw recordings through the old and current preprocessing with weights and fitted state fixed; a change isolates the processing path, which is where the schema change would show. Row 3: compare old- and new-microphone cohorts through one validated pipeline and one rubric, accounting for differences in speakers and rooms; a remaining gap points to hardware coverage. Step 3: the strip is the chapter's figure 18, the general order for any regression: drift over time, a training and validation gap, a validation and production gap, then subgroups. A yes at any check is a candidate cause, not the end; run all four, because faults often occur together. Step 4: each row of the table needs a field of the release from the retrieval slide; without the release, none can be run. WARN: students jump to ``retrain on new data.'' Retraining before diagnosis can bake a rubric change into the model. FLEX: [CORE].}
\end{frame}

\begin{frame}{Data debt has four kinds, each with a number to watch}
  \centering
  \footnotesize
  \renewcommand{\arraystretch}{1.25}
  \begin{tabular}{@{}lP{0.36\textwidth}ll@{}}
    \toprule
    \textbf{Debt} & \textbf{In the KWS regression} & \textbf{Metric} & \textbf{Warn / critical} \\
    \midrule
    Freshness & New microphone; training set now stale & PSI vs.\ baseline & 0.1 / 0.25 \\
    Quality & Labels follow a revised rubric & Known label error rate & 1\% / 5\% \\
    Schema & Device metadata changed silently & Schema version branches & 3 / 10 \\
    Documentation & Nobody recorded what changed & Datasets with data cards & 80\% / 50\% \\
    \bottomrule
  \end{tabular}\\[0.45cm]
  \uncover<2->{\normalsize A recorded release makes every row measurable.}
  \mlsyscredit{Data Engineering, \S4.12.2 and table 9 (example thresholds, not standards).}
  \note{[2 min] NARRATE: data debt is the accumulated cost of hidden coupling and missing records across the data stack; unlike a code defect, it lets every computation finish while the outputs quietly get worse. The KWS regression contains all four kinds. Freshness: the new microphone left the training set stale. Quality: labels followed a revised rubric never reconciled with the old evaluation. Schema and documentation: the device metadata changed and nobody recorded it. Read one row across: each kind has a metric and example warning and critical thresholds, so debt becomes a service-level objective a team can track. Note that the freshness metric is the PSI we just computed. WARN: these thresholds are the chapter's examples, not industry standards; calibrate them. NARRATE: the chapter models the remediation backlog as growing geometrically, $\text{Debt}_0(1+r)^n$, a planning heuristic: small differences in growth rate become large backlogs. Deliberate debt to hit a launch is a choice; unmeasured debt is the failure. Landing line: a regression with no model change is debugged by comparing fixed versions, which only a recorded release makes possible. FLEX: [OPTIONAL].}
\end{frame}

% =============================================================================
% CLOSE
% =============================================================================
\begin{frame}{Each silent failure had a number that would have shown it}
  \centering
  \scalebox{0.75}{\hookstall[0]}\hfill\scalebox{0.75}{\hookdrift[0]}\\[0.35cm]
  \begin{minipage}[t]{0.46\textwidth}\centering\footnotesize
    \uncover<1->{\textcolor{mlsysaccent}{\textbf{Watch:}} waiting share of each step\\against equation 4.\\\textcolor{mlsysaccent}{\textbf{Fix:}} the binding stage: tier, codec, shards, workers.}
  \end{minipage}\hfill
  \begin{minipage}[t]{0.46\textwidth}\centering\footnotesize
    \uncover<2->{\textcolor{mlsysaccent}{\textbf{Watch:}} K-S and PSI on per-band statistics.\\\textcolor{mlsysaccent}{\textbf{Trace:}} three fixed-version comparisons\\against the release.}
  \end{minipage}\\[0.3cm]
  \uncover<3->{\small Accelerator utilization fell from 90\% to 50\% after a data refresh (hypothetical).\\Which two numbers do you check first?}
  \mlsyscredit{Left: Mohan et al., VLDB 2021. Right: Data Engineering, \S4.12 (illustrative scenario).}
  \note{[3 min] Synthesis: return to the hook. Step 1: the stalled training runs violated the min-of-rates law. The number that exposes them is the share of each training step spent waiting on input, set against compute time; the arithmetic of equation 4 then names the stage to resize, whether storage tier, codec, shard layout, or loader workers. Step 2: the keyword spotter's regression showed first as input drift, which devices can report as per-band energy statistics compared by K-S and PSI, and was traced by three controlled comparisons against the release Part 1 built. Step 3: ENGAGE: the integrating question, one minute. Expected: bytes per sample after the refresh (did the format or image size change?) and the supply rate of the storage tier against 480 MB/s; then the waiting share per step. A second good answer: whether decode moved back onto the CPU. WARN: students who answer ``drift'' have mixed the two halves: utilization is a delivery symptom; drift changes accuracy, not utilization. FLEX: [CORE].}
\end{frame}

\begin{frame}{Takeaways}
  \normalsize
  \begin{itemize}\setlength{\itemsep}{0.3cm}
    \item Moving a petabyte costs about a day and \$20,000, so large datasets pull the computation to them.
    \item Training runs at the slower of compute and data supply.
    \item Bytes per sample is a lever as strong as bandwidth; which wins depends on the binding stage.
    \item Drift statistics flag; controlled comparisons against a release diagnose.
  \end{itemize}
  \note{[1.5 min] NARRATE: one line per pair of segments. Ask students to write the one they would explain to a teammate whose accelerator utilization dropped. FLEX: [CORE].}
\end{frame}

\begin{frame}{Open questions}
  \normalsize
  \begin{itemize}\setlength{\itemsep}{0.35cm}
    \item The ResNet-50 pipeline passes on a SATA drive today. How much headroom would you buy, and what evidence would tell you the rate is about to rise?
    \item Labels arrive weeks after a drift alert. How long should a team wait before acting, and who bears the cost of waiting?
    \item Moving compute to the data avoids egress fees but ties a team to one provider's region. When is moving a petabyte the right call?
  \end{itemize}
  \note{[1 min] No single answer. Q1: the chapter sizes storage for the rate the pipeline is heading toward; evidence includes faster kernels on the roadmap, a new accelerator, or a waiting share that grows with each release. Q2: waiting costs accuracy for the users who meet the drifted model; acting early risks retraining on a shift that did not matter; Lectures 14 and 14b develop the retraining decision. Q3: a long-lived dataset that many teams will train on repeatedly, or a region with compute the data's region lacks. FLEX: [OPTIONAL]; assign one as a short written reflection.}
\end{frame}

\begin{frame}{Read and explore}
  \small
  \begin{itemize}\setlength{\itemsep}{0.2cm}
    \item \textbf{Read:} Data Engineering, \S4.9 to \S4.12; Napkin Math 4.5, 4.7, 4.8, 4.9; the self-check questions.
    \item Mohan et al., ``Analyzing and Mitigating Data Stalls in DNN Training,'' VLDB 2021.
    \item Leclerc et al., ``FFCV: Accelerating Training by Removing Data Bottlenecks,'' CVPR 2023.
    \item Recht et al., ``Do ImageNet Classifiers Generalize to ImageNet?'' ICML 2019.
    \item Sculley et al., ``Hidden Technical Debt in Machine Learning Systems,'' NeurIPS 2015.
  \end{itemize}
  \note{[1 min] The chapter first: every number in today's delivery and drift arguments is there. The four papers each back one segment: measured data stalls and a loader that removes them (Mohan), what redesigning format and loader together buys (Leclerc), what a small shift in collection costs (Recht), and the origin of the data-debt vocabulary (Sculley). FLEX: [OPTIONAL].}
\end{frame}

\begin{frame}[c]
  \centering
  {\small\textcolor{lightgray}{NEXT LECTURE \textbar{} NEURAL COMPUTATION}}\\[0.5cm]
  {\LARGE\bfseries What does a neural network\\actually compute?}
  \note{[0.5 min] Close on the question Lecture 5 answers. Today the model was a consumer: it ate 24.6 GFLOP of work per training image, and we only asked whether the data arrived fast enough. Next time we open it and count those operations. Ask students to read Neural Computation before then.}
\end{frame}

\end{document}
